\documentclass[11pt]{article}
\usepackage[margin=1in]{geometry}
\usepackage{amsmath,amssymb,amsthm}
\usepackage{algorithm}
\usepackage{algpseudocode}
\usepackage{booktabs}
\usepackage{enumitem}
\usepackage[hidelinks]{hyperref}

\newcommand{\R}{\mathbb{R}}
\newcommand{\E}{\mathbb{E}}
\DeclareMathOperator*{\argmax}{arg\,max}
\DeclareMathOperator*{\argmin}{arg\,min}

\title{\textbf{geneal}: Active Learning for Risk-Aware Knockout Target Nomination\\
\large A Step-by-Step Method Tutorial}
\author{}
\date{}

\begin{document}
\maketitle

\begin{abstract}
This document explains, from the ground up, an active-learning method for choosing
which gene knockouts to test next so that we efficiently find targets that are
\emph{effective} (kill the cancer cell), \emph{safe} (do not kill normal cells), and
\emph{mechanistically diverse} (spread across biological pathways so we are not betting
everything on one mechanism). It is written for a reader new to active learning: every
equation is followed by a plain-language reading of what each symbol and term means and
how to interpret it. Data collection and gene representation are out of scope.
\end{abstract}

\tableofcontents

\section{What problem are we solving, and what is active learning?}

We have a large list of candidate genes. For a given cancer cell line, knocking out a
gene has some effect on the cell's survival; a strongly negative effect means the cell
dies, which is what we want for a cancer target. Measuring this for every gene is
expensive, so we cannot test them all. We want to test a small number, in batches, and
learn as we go which genes are worth testing next.

\paragraph{Active learning (AL), in one sentence.} Instead of choosing what to measure
all at once or at random, we build a predictive model from the measurements we already
have, use it to decide which \emph{unmeasured} genes are most worth measuring next,
measure those, update the model, and repeat. The aim is to reach a good answer with as
few measurements as possible.

\paragraph{The loop.} One ``round'' of active learning is:
\begin{enumerate}[leftmargin=2em,itemsep=1pt]
\item Fit a model on the genes measured so far.
\item Score every not-yet-measured gene by how useful it would be to measure (the
\emph{acquisition function}).
\item Pick the top batch of $q$ genes, measure them, add them to the measured set.
\item Repeat.
\end{enumerate}

\paragraph{Notation we will use throughout.}
\begin{itemize}[leftmargin=2em,itemsep=1pt]
\item $\mathcal{G}=\{1,\dots,N\}$: the candidate genes. $N$ is the total number.
\item $x_i \in \R^d$: a fixed numeric ``fingerprint'' (embedding) of gene $i$, a vector of
$d$ numbers. Genes with similar biology have similar $x_i$.
\item $e_i$: the \emph{efficacy} of knocking out gene $i$ $=$ how lethal it is to the
cancer cell. Higher $=$ better target.
\item $t_i$: the \emph{toxicity} $=$ how lethal the knockout is in contexts we want to
spare (normal-like cells). Higher $=$ worse / more dangerous.
\item $D \subset \mathcal{G}$: the set of genes already measured (``revealed'').
\item $B \subset \mathcal{G}$, $|B|=K$: the final shortlist (``portfolio'') of nominated targets.
\end{itemize}
Everything random uses a fixed seed, so a run can be reproduced exactly, and any two
methods we compare start from the same initial measurements (this is called using
\emph{common random numbers}, and it makes the comparison fair).

\section{Step 1: the surrogate model (a Gaussian process)}

\paragraph{Why we need it.} To decide what to measure next we need a model that, from the
genes measured so far, predicts the efficacy of any unmeasured gene \emph{and} tells us
how confident it is. That confidence is essential: a gene we are unsure about may be
worth measuring precisely because measuring it teaches us a lot. A Gaussian process (GP)
is a model that outputs exactly this: for any gene, a predicted value \emph{and} an
uncertainty.

\paragraph{The idea.} A GP assumes that genes with similar fingerprints $x$ tend to have
similar efficacy, and it quantifies ``similar'' with a \emph{kernel} function $k(x,x')$ ---
a similarity score between two genes. We write
\begin{equation}
y(x)\sim\mathcal{GP}\big(m(x),\,k(x,x')\big).
\end{equation}
\textbf{How to read this:} ``the unknown function $y(\cdot)$ mapping a gene's fingerprint
to its efficacy is modeled as a Gaussian process with a baseline level $m(x)$ (a constant
we fit) and a similarity rule $k(x,x')$.'' The GP is a distribution over \emph{functions}:
before seeing data it represents all plausible efficacy-vs-fingerprint curves; after seeing
data it narrows to those consistent with the measurements.

\paragraph{The kernel (similarity rule).} We use the Mat\'ern-$5/2$ kernel:
\begin{equation}
k(x,x') = \theta\Big(1+\tfrac{\sqrt{5}\,r}{\ell}+\tfrac{5\,r^2}{3\ell^2}\Big)\exp\!\Big(-\tfrac{\sqrt5\, r}{\ell}\Big),\qquad r=\lVert x-x'\rVert .
\end{equation}
\textbf{Reading it term by term:}
\begin{itemize}[leftmargin=2em,itemsep=1pt]
\item $r=\lVert x-x'\rVert$ is the distance between the two genes' fingerprints. Small $r$
$=$ similar genes.
\item The $\exp(-\sqrt5\,r/\ell)$ factor makes similarity \emph{decay with distance}: far-apart
genes are nearly unrelated. So a measurement of one gene informs predictions mostly for
nearby genes.
\item $\ell$ (the ``length scale'') sets \emph{how fast} that decay happens --- how far the
influence of a measurement reaches. Large $\ell$ $=$ smooth, long-range sharing; small
$\ell$ $=$ only very close genes share information.
\item $\theta$ scales the overall variance --- how much efficacy varies across genes.
\item The polynomial $(1+\dots)$ makes the function ``mildly smooth'' (Mat\'ern-$5/2$): smoother
than a jagged model, but not as unrealistically smooth as the common ``RBF'' kernel.
\end{itemize}
We learn $\theta,\ell$ (and a noise level, below) automatically from the data.

\paragraph{What the GP gives you: prediction $+$ uncertainty.} After measuring genes
$X_D$ with (noisy) efficacy values $\mathbf{y}_D$, the GP's prediction for a new gene
$x_\ast$ is a Gaussian with mean $\mu(x_\ast)$ and variance $\sigma^2(x_\ast)$:
\begin{align}
\mu(x_\ast) &= k_\ast^\top (K+\sigma_n^2 I)^{-1}\,\mathbf{y}_D, \label{eq:gpmean}\\
\sigma^2(x_\ast) &= \underbrace{k(x_\ast,x_\ast)}_{\text{prior uncertainty}} - \underbrace{k_\ast^\top (K+\sigma_n^2 I)^{-1} k_\ast}_{\text{uncertainty removed by data}}. \label{eq:gpvar}
\end{align}
\textbf{How to read \eqref{eq:gpmean} (the prediction):}
\begin{itemize}[leftmargin=2em,itemsep=1pt]
\item $k_\ast$ is the vector of similarities between the new gene $x_\ast$ and each measured
gene. $(k_\ast)_i=k(x_i,x_\ast)$.
\item $K$ is the matrix of similarities \emph{among} the measured genes, $K_{ij}=k(x_i,x_j)$.
\item $\sigma_n^2$ is the measurement-noise level (assays are imperfect); $I$ is the identity.
\item In words: the prediction is a \emph{weighted average of the measured efficacies}
$\mathbf{y}_D$, where measured genes more similar to $x_\ast$ get more weight (through
$k_\ast$), after correcting for redundancy among the measured genes (the $(K+\sigma_n^2 I)^{-1}$).
\end{itemize}
\textbf{How to read \eqref{eq:gpvar} (the uncertainty):} it starts at the prior uncertainty
$k(x_\ast,x_\ast)$ and \emph{subtracts} how much the existing measurements pin down $x_\ast$.
So the variance is small (we are confident) for genes near where we have measured, and
large (we are unsure) for genes in unexplored regions. This is the quantity active learning
exploits to decide where to explore.

\paragraph{Fitting the kernel.} We choose $(\theta,\ell,\sigma_n)$ to maximize the
\emph{log marginal likelihood},
\begin{equation}
\log p(\mathbf{y}_D\mid X_D)= \underbrace{-\tfrac12 \mathbf{y}_D^\top(K+\sigma_n^2 I)^{-1}\mathbf{y}_D}_{\text{fits the data}}
\ \underbrace{-\tfrac12\log\lvert K+\sigma_n^2 I\rvert}_{\text{penalizes complexity}}\ -\tfrac{n}{2}\log 2\pi .
\end{equation}
\textbf{Reading it:} the first term rewards kernels that explain the measured values well;
the second term penalizes overly flexible kernels (Occam's razor, automatically balancing
fit against simplicity). The last term is a constant. Maximizing this picks the
similarity rule best supported by the data. When we later need two unknowns (efficacy and
toxicity) we simply fit one independent GP for each.

\section{Step 2: scoring genes to measure (single-objective acquisition)}

The \emph{acquisition function} turns the GP's prediction and uncertainty into a single
``how worth measuring is this gene'' score. Two standard choices:

\paragraph{Upper Confidence Bound (UCB).}
\begin{equation}
a_{\mathrm{UCB}}(x)=\underbrace{\mu(x)}_{\text{exploit}}+\beta\underbrace{\sigma(x)}_{\text{explore}} .
\end{equation}
\textbf{Reading it:} score a gene by its predicted efficacy $\mu$ \emph{plus} a bonus for
uncertainty $\sigma$. The first term ``exploits'' (favors genes we think are good); the
second ``explores'' (favors genes we are unsure about, which could be good and would teach
us a lot). $\beta\ge 0$ is the knob that sets the balance: $\beta=0$ is pure greed,
large $\beta$ is aggressive exploration. We maximize this score.

\paragraph{Expected Improvement (EI).}
\begin{equation}
a_{\mathrm{EI}}(x)=(\mu(x)-y^+)\,\Phi(z)+\sigma(x)\,\phi(z),\qquad z=\frac{\mu(x)-y^+}{\sigma(x)} .
\end{equation}
\textbf{Reading it:} $y^+$ is the best efficacy we have measured so far (the ``incumbent'').
$z$ measures how many standard deviations the gene's predicted value sits above that
incumbent. $\Phi$ and $\phi$ are the standard normal CDF and PDF. The whole expression is
the \emph{expected amount by which this gene would beat the current best}, accounting for
uncertainty: it is large either when the mean clearly beats $y^+$, or when the gene is
uncertain enough to plausibly beat it. Higher $=$ more worth measuring.

\section{Step 3: choosing a non-redundant batch}

We measure $q>1$ genes per round. Naively taking the top-$q$ by acquisition score can give
$q$ near-duplicate genes (same pathway, near-identical fingerprints) --- wasteful, because
once you have measured one you learn little extra from its twins. We want a batch that is
high quality \emph{and} varied.

\subsection{Determinantal point process ($k$-DPP) --- diversity from a determinant}

A determinantal point process gives each candidate subset $B$ a ``score'' equal to the
determinant of a small matrix built from $B$:
\begin{equation}
\text{score}(B)\;=\;\det\!\big(L_B\big),\qquad L_{ij}=\underbrace{q_i}_{\text{quality of }i}\ \underbrace{S_{ij}}_{\text{similarity }i,j}\ \underbrace{q_j}_{\text{quality of }j} .
\end{equation}
\textbf{Reading the pieces:}
\begin{itemize}[leftmargin=2em,itemsep=1pt]
\item $q_i$ is the \emph{quality} of gene $i$ --- e.g.\ its UCB score. It appears on both sides,
so the diagonal of $L$ is $q_i^2$: high-quality genes have large diagonal entries.
\item $S_{ij}\in[-1,1]$ is the \emph{similarity} between genes $i$ and $j$ in outcome space,
$S_{ij}=\Sigma_{ij}/\sqrt{\Sigma_{ii}\Sigma_{jj}}$, the correlation from the GP's joint
posterior covariance $\Sigma$. $S_{ij}\approx 1$ means ``measuring one tells you the
other''; $S_{ij}\approx 0$ means they are independent.
\item $\det(L_B)$ is large only when the chosen genes are \emph{both} high quality (big
diagonal) \emph{and} dissimilar (small off-diagonal, so the rows are nearly perpendicular).
Two near-identical high-quality genes make two nearly-parallel rows, which collapses the
determinant toward zero --- so the DPP will not pick both.
\end{itemize}
\textbf{Why this is principled (the key fact).} For Gaussian variables, $\tfrac12\log\det$
of a covariance is the \emph{entropy} (information content). So maximizing $\log\det(L_B)$
literally maximizes the information the batch carries: it is at once a quality objective
and a diversity objective, not an ad-hoc combination. Finding the exact best size-$q$ set
is computationally hard, so we add genes one at a time, each time taking the one that
increases $\log\det$ the most (a fast ``greedy'' rule).

\subsection{A simpler, robust hedge: cap genes per pathway}

At genome scale the similarity $S$ is very sparse (almost no two genes are correlated), so
the DPP's diversity term barely fires and a simpler tool is preferable: directly limit how
many chosen genes may come from the same pathway. Given a map $\pi(i)=$ ``the set of
pathways gene $i$ belongs to'' and a cap $c$, we go down the quality-ranked list and accept
a gene only if none of its pathways has already hit $c$ members. \textbf{Why:} this bounds
how concentrated the portfolio is in any single pathway using one interpretable number
$c$, which matters for risk (Step 6), and it does not rely on the (often absent) similarity
structure.

\section{Step 4: defining toxicity and the safety constraint}

\paragraph{Three ways to define toxicity.}
\begin{enumerate}[leftmargin=2em,itemsep=2pt]
\item \textbf{Common-essential (known in advance).}
$c_i=\tfrac{1}{L}\sum_{\ell=1}^{L}\mathbf{1}[\mathrm{effect}_{i,\ell}<\tau_e]$.
\textbf{Reading it:} count the fraction of all $L$ cell lines in which knocking out gene
$i$ is strongly lethal ($\mathbf{1}[\cdot]$ is $1$ when true, else $0$; $\tau_e$ is the
``strongly lethal'' threshold). A gene lethal almost everywhere is ``pan-essential'' ---
it would likely kill normal cells too, so high $c_i$ $=$ toxic. This is known from prior
screens, so we treat it as given.
\item \textbf{Contrast line.} $t_i=-\,\mathrm{effect}_{i,B}$: lethality in one fixed
reference cell line $B$ chosen to stand in for normal tissue. We do \emph{not} know this in
advance; a second GP learns it.
\item \textbf{Population.} $t_i=-\tfrac1L\sum_\ell \mathrm{effect}_{i,\ell}$: average
lethality across the whole panel. Also learned.
\end{enumerate}

\paragraph{Why we use a hard constraint, not a penalty.} A tempting approach is to
score genes by $s_i=e_i-\lambda t_i$ (efficacy minus a toxicity penalty). But this treats
safety as something you can ``buy back'' with enough efficacy, and the weight $\lambda$ has
no natural meaning or units. In drug discovery, a too-toxic target is simply \emph{ruled
out}. So we impose a ceiling instead.

\section{Step 5: constrained nomination (the selection rule)}

\begin{equation}
B^\star=\argmax_{|B|=K}\ \sum_{i\in B}\hat e_i \qquad\text{subject to}\qquad \hat t_i\le\tau\ \ \text{for every } i\in B .
\label{eq:constrained}
\end{equation}
\textbf{Reading it:} ``choose the $K$ genes with the highest predicted efficacy
$\hat e_i$, but only among genes whose predicted toxicity $\hat t_i$ is at or below the
safety ceiling $\tau$.'' Concretely: throw out every gene with $\hat t_i>\tau$, then take
the top $K$ by predicted efficacy from what remains.
\begin{itemize}[leftmargin=2em,itemsep=1pt]
\item $\tau$ is the single, interpretable knob: ``the maximum toxicity I will tolerate.''
\item Sweeping $\tau$ from strict to lenient traces an \emph{efficacy--safety frontier}: a
tighter ceiling gives safer but typically less potent shortlists.
\item If toxicity is known we use $c_i$ in the constraint; if learned we use the toxicity
GP's prediction (optionally its \emph{upper} confidence bound, so we only admit genes we
are \emph{confident} are safe).
\item We then layer the per-pathway cap (Step 3.2) on top, so the shortlist is also
mechanism-diversified.
\end{itemize}

\section{Step 6: when both efficacy and toxicity are unknown (multi-objective AL)}

If we must \emph{learn} toxicity too, each measurement reveals a gene's efficacy and its
toxicity together, and we want to explore both at once. Now each gene is a \emph{point in
two dimensions} $f(x)=(e(x),\,-t(x))$ --- we write $-t$ so that ``bigger is better'' on
both axes (high efficacy, low toxicity).

\paragraph{Domination and the Pareto front.} Point $u$ ``dominates'' $v$ if $u$ is at least
as good on both axes and strictly better on one. The genes that are \emph{not} dominated by
any other form the \emph{Pareto front} --- the set of best available trade-offs (you cannot
improve efficacy without worsening toxicity, or vice versa). There is no single best point;
the front is the answer.

\paragraph{Hypervolume: scoring a whole front with one number.}
\begin{equation}
\mathrm{HV}(P;r)=\text{area}\Big(\bigcup_{p\in P}[r_1,p_1]\times[r_2,p_2]\Big).
\end{equation}
\textbf{Reading it:} pick a ``reference point'' $r$ that is worse than everything (a lower
bound on both axes). For each point $p$ in the set $P$, draw the rectangle from $r$ up to
$p$. The hypervolume is the total area covered by the union of those rectangles --- the
amount of objective space that $P$ ``dominates.'' A larger hypervolume means the set covers
more of the desirable high-efficacy/low-toxicity region. It is the standard single-number
quality score for a set of trade-offs, and we track it over rounds to see how fast each
method discovers good targets.

\paragraph{Expected Hypervolume Improvement (EHVI): what to measure next.}
\begin{equation}
a_{\mathrm{EHVI}}(x)=\E_{f(x)}\big[\,\mathrm{HV}(P\cup\{f(x)\};r)-\mathrm{HV}(P;r)\,\big]\ \ge 0 .
\label{eq:ehvi}
\end{equation}
\textbf{Reading it:} $P$ is the current front. $f(x)$ is the (uncertain, Gaussian) location
the new gene would occupy. The bracket is ``how much would the hypervolume \emph{grow} if
this gene landed where we think it might.'' The $\E[\cdot]$ averages that growth over the
GP's uncertainty about $f(x)$. So EHVI rewards genes \emph{expected to enlarge the set of
good trade-offs the most}. It is the two-dimensional version of Expected Improvement, and
it automatically balances exploit (promising mean location) against explore (large
uncertainty $\Rightarrow$ chance of a big improvement), \emph{without} needing to pick a
trade-off weight. It is never negative because a gene can only add area, never remove it.

\paragraph{How we compute it.} The expectation has no simple closed form we rely on, so we
estimate it by \emph{Monte Carlo}: draw many possible $(e,-t)$ values from the gene's GP
posterior, compute the hypervolume increase for each draw, and average:
\begin{equation}
\hat a_{\mathrm{EHVI}}(x)=\frac1M\sum_{m=1}^{M}\max\!\big(0,\ \mathrm{HV}(P\cup\{f^{(m)}\};r)-\mathrm{HV}(P;r)\big),\quad f^{(m)}\sim\mathcal{N}\!\big(\boldsymbol\mu(x),\operatorname{diag}\boldsymbol\sigma^2(x)\big).
\end{equation}
\textbf{Reading it:} $f^{(m)}$ is the $m$-th random ``what if this gene turns out here''
sample; we average the area it would add over $M$ samples. To pick a batch of $q$, we take
the best gene, pretend it landed at its predicted mean (a ``fantasy''), add that to the
front, and recompute --- so the next pick complements the first. For speed on large gene
lists we run this only on a shortlist of the most promising candidates (those with high
optimistic efficacy and low optimistic toxicity), which never drops a genuinely
front-improving gene.

\section{Step 7: judging the portfolio's risk}

Even a high-efficacy, low-toxicity shortlist can be \emph{fragile}: if all its targets sit
in one pathway and that pathway later proves toxic or undruggable, every target fails at
once. We measure this. With pathway membership $\pi$ and a per-gene value $v_i$ (e.g.\
efficacy):
\begin{align}
\text{concentration}\quad & C(B)=\max_{p}\ \frac{\#\{\text{genes in } B \text{ that belong to pathway } p\}}{\#\{\text{genes in } B \text{ with any pathway}\}}, \label{eq:conc}\\[2pt]
\text{dropout robustness}\quad & R(B)=1-\frac{1}{\#\text{pathways in }B}\sum_{p}\frac{\text{value of }B\text{'s genes in pathway }p}{\text{total value of }B}. \label{eq:rob}
\end{align}
\textbf{Reading \eqref{eq:conc}:} the single largest share of the portfolio sitting in any
one pathway. $C=1$ means everything is in one pathway (maximum risk); a small $C$ means the
bets are spread out. \emph{Lower is safer.}\\
\textbf{Reading \eqref{eq:rob}:} imagine deleting one pathway at random (it turned out
toxic), which removes all of $B$'s genes in it. The fraction summed is the share of value
lost; averaging over which pathway gets deleted and subtracting from $1$ gives the
\emph{expected fraction of value that survives}. \emph{Higher is safer.} The per-pathway
cap improves both: it lowers $C$ and raises $R$, at a measurable cost in efficacy --- this
trade-off is exactly the efficacy--risk frontier we report.

\section{Step 8: how we measure success}
As the AL loop runs we track, versus the number of rounds:
\begin{itemize}[leftmargin=2em,itemsep=2pt]
\item \textbf{recall@$k$}: of the truly $k$ most-lethal genes, what fraction have we
measured so far? (Did we find the important ones, and how fast?)
\item \textbf{max value}: the single best true efficacy among measured genes, divided by
the global best (did we at least find one excellent hit, even if the average drops?).
\item \textbf{diversity}: average pairwise fingerprint distance of everything measured ---
how spread out our exploration is.
\item \textbf{$\alpha$-NDCG@$k$}: a coverage score that rewards finding lethal genes across
\emph{distinct} mechanisms and discounts redundant ones. It groups genes into ``nuggets''
(here, clusters of similar genes); the credit for the gene at rank $i$ that covers nugget
$n$ is multiplied by $(1-\alpha)^{r}$, where $r$ counts how many earlier-ranked genes
already covered $n$. So the second, third, \dots\ gene from the same mechanism earns
geometrically less. The total is normalized against an ideal ordering so it lands in
$[0,1]$. High $=$ many lethal genes spanning many mechanisms.
\end{itemize}

\section{Putting it together}
\begin{algorithm}[h]
\caption{One round of risk-aware nomination (constrained variant)}
\begin{algorithmic}[1]
\State Fit the efficacy GP on the measured genes; if toxicity is unknown, fit a toxicity GP too.
\State For every unmeasured gene, predict efficacy $\hat e_i$ and toxicity $\hat t_i$ (with uncertainties).
\State Keep only genes that clear the safety ceiling: $\mathcal{E}=\{i:\hat t_i\le\tau\}$.
\State Build the shortlist $B$: take the highest-$\hat e_i$ genes in $\mathcal{E}$, allowing at most $c$ per pathway.
\State (Exploration) decide which genes to \emph{measure} next via UCB / $k$-DPP / EHVI; measure them; add to the measured set.
\State Report efficacy, toxicity, pathway concentration and dropout robustness, and recall@$k$.
\end{algorithmic}
\end{algorithm}

In short: the \textbf{GP} gives calibrated predictions and uncertainty; the
\textbf{acquisition functions} (UCB / $k$-DPP / EHVI) decide what to measure to learn
fastest; the \textbf{constraint} enforces safety as a hard ceiling; the \textbf{per-pathway
cap} hedges against pathway-level failure; and the \textbf{metrics} quantify the
efficacy--safety--diversity trade-offs a target-nomination pipeline must navigate.

\begin{thebibliography}{9}
\bibitem{rasmussen} C.~E.~Rasmussen and C.~K.~I.~Williams. \emph{Gaussian Processes for Machine Learning}. MIT Press, 2006.
\bibitem{settles} B.~Settles. \emph{Active Learning Literature Survey}. Tech.\ Report, 2009.
\bibitem{kulesza} A.~Kulesza and B.~Taskar. \emph{Determinantal Point Processes for Machine Learning}. Found.\ Trends ML, 2012.
\bibitem{emmerich} M.~Emmerich et al. Single- and multiobjective optimization assisted by Gaussian random field metamodels. \emph{IEEE TEC}, 2006.
\bibitem{clarke} C.~L.~A.~Clarke et al. Novelty and diversity in information retrieval evaluation. \emph{SIGIR}, 2008.
\end{thebibliography}

\end{document}
